\section{FAQ 与经验总结}

\subsection{常见问题}

\begin{itemize}
    \item \textbf{Q: DSPy 适合多大项目？} A: 适合所有需要组合多个 LLM 组件的项目，尤其是当你希望用代码治理 prompt 时。
    \item \textbf{Q: Evidence Matrix 需要保留多久？} A: 至少保留 30 天或上一轮 rollout 的时间窗口，便于做 regress reviews。
    \item \textbf{Q: 怎么处理 optimizer drift？} A: 先回溯 Evidence Matrix，再用 fallback prompt 进行对比跑，同时记录 metric delta。
    \item \textbf{Q: 是否需要专门的 DSL？} A: DSPy 自带 DSL，只需在 Python 中声明模块和 signature，无需自己造轮子。
    \item \textbf{Q: 如何衡量 optimizer improvements？} A: 以 `automation coverage`, `latency`, `metric accuracy` 三大指标组合来量化。
\end{itemize}

\subsection{经验提炼}

经验告诉我们：1) 每次 release 要把 evidence snapshot 交给 QA；2) monitor 需要 24/7 alerting + weekly review；3) documentation 不能留在脑子里，必须写成 `DSPy.md` + `Incident Log`。

\subsection{本章小结}

FAQ 与经验总结让团队对 DSPy pipeline 的常见问题有快速解答，避免在实践中重复踩坑。

